\subsection{Vision Transformer}
\label{sec:vision_transformer}

\begin{figure}[htbp]
\centering
\includegraphics[width=\columnwidth]{vit.png}
\caption{Architecture of the ViT Backbone and Patch Projection.}
\label{fig:3}
\end{figure}

The visual representation extraction architecture employs the Vision Transformer (ViT) Base backbone, processing two-dimensional facial images as sequences of one-dimensional patch tokens, as illustrated in Figure~\ref{fig:3} \cite{ref26}. An input image $\mathbf{x} \in \mathbb{R}^{H \times W \times C}$ with resolution $224 \times 224 \times 3$ is partitioned into $N = 196$ non-overlapping patches of size $P \times P = 16 \times 16$ pixels \cite{ref27}. Each patch is linearly projected via matrix $\mathbf{E} \in \mathbb{R}^{(P^2 C) \times D}$ into a latent embedding space of dimension $D = 768$, and concatenated with a class token $\mathbf{x}_{\text{class}}$ and positional embeddings $\mathbf{E}_{\text{pos}}$ as formulated in \eqref{eq:1} \cite{ref27}. The token sequence propagates through $L = 12$ transformer encoder layers, each comprising a Multi-Head Self-Attention (MHSA) module and an MLP block with Layer Normalization (LN) as in \eqref{eq:2} and \eqref{eq:3} \cite{ref28}. The self-attention mechanism captures global spatial dependencies across all image patches without receptive field constraints inherent to conventional convolutions \cite{ref29}.

\begin{equation}
\mathbf{z}_0 = [\mathbf{x}_{\text{class}}; \, \mathbf{x}_p^1\mathbf{E}; \, \dots; \, \mathbf{x}_p^{196}\mathbf{E}] + \mathbf{E}_{\text{pos}}
\label{eq:1}
\end{equation}

\begin{equation}
\mathbf{z}'_\ell = \text{MHSA}(\text{LN}(\mathbf{z}_{\ell-1})) + \mathbf{z}_{\ell-1}
\label{eq:2}
\end{equation}

\begin{equation}
\mathbf{z}_\ell = \text{MLP}(\text{LN}(\mathbf{z}'_\ell)) + \mathbf{z}'_\ell
\label{eq:3}
\end{equation}

\begin{table}[htbp]
\caption{Multi-Domain Feature Fusion and Ablation Configurations.}
\label{tab:II}
\centering
\begin{tabular}{clcc}
\toprule
\# & Configuration & Domain Category & Dimension \\
\midrule
1 & Face & Single-Domain & 768 \\
2 & Emotion & Single-Domain & 768 \\
3 & Age & Single-Domain & 768 \\
4 & Emotion $\oplus$ Face & Dual-Domain & 1,536 \\
5 & Face $\oplus$ Age & Dual-Domain & 1,536 \\
6 & Emotion $\oplus$ Age & Dual-Domain & 1,536 \\
7 & Face $\oplus$ Emotion $\oplus$ Age & Tri-Domain (Proposed) & 2,304 \\
\bottomrule
\end{tabular}
\end{table}

To ensure computational efficiency, the three ViT backbones remain frozen as offline feature extractors. The checkpoints comprise ViT-Face (`skutaada/\allowbreak VIT-\allowbreak VGGFace`) capturing facial biometric geometry, ViT-Emotion (`dima806/\allowbreak facial\_\allowbreak emotions\_\allowbreak image\_\allowbreak detection`) capturing affective dynamic expressions, and ViT-Age (`dima806/\allowbreak facial\_\allowbreak age\_\allowbreak image\_\allowbreak detection`) capturing biological aging morphology. For each model, latent feature representation $\mathbf{f}_{\text{domain}}$ is extracted from the [CLS] token of final encoder layer $L$ via Layer Normalization in \eqref{eq:4}. Multi-domain feature fusion integrates extracted latent vectors via direct sum concatenation, producing a 2,304-dimensional tri-domain vector $\mathbf{z}_{\text{tri}}$ in \eqref{eq:5}. To systematically evaluate individual and joint domain contributions, seven feature representation ablation configurations are formulated as detailed in Table~\ref{tab:II}.

\begin{equation}
\mathbf{f}_{\text{domain}} = \text{LN}(\mathbf{z}_L^0) \in \mathbb{R}^{768}
\label{eq:4}
\end{equation}

\begin{equation}
\mathbf{z}_{\text{tri}} = \mathbf{f}_{\text{face}} \oplus \mathbf{f}_{\text{emotion}} \oplus \mathbf{f}_{\text{age}} \in \mathbb{R}^{2304}
\label{eq:5}
\end{equation}
